% TOP_PROBLEM: {"id": 3252, "title": "List Colourings of Complete Multipartite Graphs with 2 Big Parts", "category_id": 3, "category": {"id": 3, "name": "graph_theory", "display_name": "Graph Theory", "description": "Problems involving graphs, networks, and their properties.", "slug": "graph-theory", "order_index": 3, "created_at": "2026-07-31T15:26:25.671Z"}, "status": "open", "problem_number": "OPG-59927", "set_id": 12}
% FIRST_POSTED: 2026-10-02
% ATTEMPT_STATUS: unresolved
% UnsolvedMath ID 3252; source catalogue code OPG-59927
% !TeX program = lualatex
% Research attempt by GPT-6 Astra Ultra; no claim of a new complete solution.
\documentclass[11pt,a4paper]{article}
\usepackage[margin=25mm]{geometry}
\usepackage{fontspec}
\setmainfont{lmroman10-regular.otf}[BoldFont=lmroman10-bold.otf,
 ItalicFont=lmroman10-italic.otf,BoldItalicFont=lmroman10-bolditalic.otf]
\usepackage{amsmath,amssymb,mathtools}
\usepackage{unicode-math}
\setmathfont{latinmodern-math.otf}
\AtBeginDocument{\let\setminus\smallsetminus}
\usepackage{xurl,hyperref}
\hypersetup{colorlinks=true,urlcolor=blue}
\setcounter{secnumdepth}{0}
\setlength{\parindent}{0pt}
\setlength{\parskip}{5pt}
\setlength{\emergencystretch}{3em}
\begin{document}
\section*{List Colourings of Complete Multipartite Graphs with 2 Big Parts}\label{3252}
{\small UnsolvedMath ID 3252 · OPG-59927 · Graph Theory · OpenGarden}
\subsection{Definitions and mathematical statement}
Fix integers $a,b\ge2$. The complete bipartite graph $K_{a,b}$ has disjoint
independent vertex sets of sizes $a$ and $b$, with every edge between the
two sets. The join $G+H$ of disjoint graphs adds every edge between
$V(G)$ and $V(H)$. Write $K_t$ for the complete graph on $t$ vertices,
including $K_0$ with no vertices. Thus $K_{a,b}+K_t$ is a complete
multipartite graph with part sizes $a,b,1,\ldots,1$.

A proper colouring assigns different colours to adjacent vertices.
The chromatic number of this join is $t+2$: the two large parts require
two colours, and each of the $t$ joined clique vertices requires its own
additional colour. Its list chromatic number $\chi_\ell$ is the least
integer $q$ such that every assignment of $q$ permitted colours to each
vertex admits a proper colouring from those lists.

Determine, for every $a,b\ge2$, the parameter
\[
 \phi(a,b)=\min\{t\in\mathbb N_0:
                 \chi_\ell(K_{a,b}+K_t)=t+2\}.
\]
The source discusses a theorem ensuring the defining set is nonempty.
A complete determination must supply both a universal list-colouring
argument at the proposed $t$ and, for every smaller $t$, a list assignment
that cannot be coloured with $t+2$ available colours per vertex. Ordinary
chromatic number alone does not decide the threshold. Interchanging $a$
and $b$ leaves the graph unchanged, so the sought function is symmetric.
\subsection{Short English statement}
How many mutually adjacent universal vertices must be added to a complete bipartite graph before its choice number equals its chromatic number?
\subsection{Research attempt}
\paragraph{Scope and process.}
This is a GPT-6 Astra Ultra research attempt dated 2 October 2026, using
primary-source browsing and elementary mathematical checks. The canonical
definition above is retained. Results below are restricted results or
reductions, not a claim of a new solution to the entire catalogue question.
No formal proof assistant or external mathematical referee checked this text.


\paragraph{Literature and approach.}
The original question [S2] attributes this threshold problem to Allagan
and records bounds stronger than the elementary ones below for some
unbalanced parameters. Noel, Reed and Wu [R1] prove Ohba's theorem:
a graph with at most $2\chi+1$ vertices is chromatic-choosable. We
use that theorem explicitly as an external input for upper bounds,
while independently proving the small obstructions, the zero-threshold
classification and monotonicity facts. These deductions do not determine
$\phi(a,b)$ for every pair.

\paragraph{The parameter really is a threshold.}
Suppose $K_{a,b}+K_t$ is $(t+2)$-choosable. Given lists of size
$t+3$ on $K_{a,b}+K_{t+1}$, colour one universal clique vertex
arbitrarily. Remove its colour from all other lists. Every remaining
list has at least $t+2$ colours, so the assumed property colours
the remaining graph. Hence the defining property, once true, stays
true for every larger $t$.

For fixed $t$, deleting vertices from either large part preserves
$(t+2)$-choosability, and the ordinary chromatic number remains $t+2$
as long as both parts are nonempty. Thus $\phi(a,b)$ is symmetric
and is nondecreasing in each of its two arguments. These facts follow
from actual list-colouring extensions and restrictions, rather than
from ordinary chromatic-number arithmetic alone.

\paragraph{A complete classification of when $t=0$ works.}
We first show that $K_{2,b}$ is two-choosable for $b\le3$. Give
the two vertices in the small part lists $L_1,L_2$ of size two.
If their lists intersect, colour both by a common colour; each vertex
on the other side retains another permitted colour. If the lists are
disjoint, there are four choices of an unordered pair consisting of
one colour from each list. A two-element list at a vertex on the
other side forbids at most one such pair: it is unusable only if
it equals precisely that chosen pair. At most three vertices forbid
at most three pairs, so one of the four choices works. Every other
vertex then has a colour outside the chosen pair, completing the
proper colouring.

For $K_{2,4}$, give the small side lists $\{1,2\},\{3,4\}$,
and give the other side the four lists $\{1,3\},\{1,4\},\{2,3\},
\{2,4\}$. Whatever pair is chosen on the small side, the vertex
with precisely that pair as its list cannot be coloured.
For $K_{3,3}$, assign the three lists $\{1,2\},\{1,3\},\{2,3\}$
once on each side. Each side must use at least two of the three
colours, so the used-colour sets intersect, impossible across a complete
bipartite join.

By taking these graphs as induced subgraphs, the obstructions extend
to every pair $a,b\ge2$ except $(2,2),(2,3),(3,2)$. Therefore
\[
 \phi(a,b)=0\quad\Longleftrightarrow\quad
 \{a,b\}=\{2,2\}\text{ or }\{2,3\}.
\]
The companion script [R2] checks the two explicit bad list assignments
by exhaustive enumeration of their permitted vertex colours.

\paragraph{Upper bounds and two exact positive thresholds.}
For $H=K_{a,b}+K_t$, its order is $a+b+t$ and its chromatic
number is $t+2$. Ohba's theorem applies whenever
\[
 a+b+t\le2(t+2)+1,
 \quad\text{equivalently }t\ge a+b-5.
\]
Hence $\phi(a,b)\le\max\{0,a+b-5\}$. At $(a,b)=(2,4)$
and $(3,3)$ this gives an upper bound one, and the preceding bad
lists give the matching lower bound. Thus
$\phi(2,4)=\phi(4,2)=\phi(3,3)=1$. The upper direction here is
credited to [R1]; it is not independently proved by the finite tests.

More generally, induced-subgraph monotonicity gives
$\phi(a,b)\ge\chi_\ell(K_{a,b})-2$. To see a concrete family,
put $n=\binom{2q-1}{q}$ and assign all $q$-subsets of a
$(2q-1)$-colour universe to each side of $K_{n,n}$. Each side
must use at least $q$ colours, giving a conflict. Thus
$\chi_\ell(K_{n,n})>q$ and $\phi(n,n)\ge q-1$. This shows
that the required number of universal vertices cannot stay bounded
as both parts grow, although it leaves a large gap to the general
upper bound.

\paragraph{Final status and first remaining gap.}
\textbf{UNRESOLVED.} The zero threshold and the three displayed
threshold-one pairs are established, with exact obstructions and a
clearly identified external upper theorem. For general $a,b$, the
missing work is a tight universal list-colouring construction together
with uncolourable assignments at every smaller $t$; neither monotonicity
nor the order theorem supplies that sharp pair-dependent analysis.

\subsection{Sources}
\begin{itemize}
\item[S1] ULAM.AI and the UnsolvedMath Contributors, supplied problems.json, record 3252 (OPG-59927), OpenGarden collection. Catalogue identity and source transcription; CC BY 4.0.
\url{https://huggingface.co/datasets/ulamai/UnsolvedMath}.
\item[S2] Open Problem Garden, List Colourings of Complete Multipartite Graphs with 2 Big Parts. Original problem and discussion; the mathematical formulation below is expanded from the supplied source record.
\url{http://www.openproblemgarden.org/op/list_colourings_of_complete_multipartite_graphs_with_2_big_parts}.

\item[R1] Jonathan A. Noel, Bruce A. Reed and Hehui Wu,
\emph{A Proof of a Conjecture of Ohba}, Journal of Graph Theory
79 (2015), 86--102. Primary theorem checked 2 October 2026.
\url{https://arxiv.org/abs/1211.1999}.
\item[R2] GPT-6 Astra Ultra, exact batch certificates, 2 October 2026.
Repository script \texttt{scripts/check\_attempt\_batch03\_colouring\_2026\_10\_02.py}.

\end{itemize}
\end{document}
